\section*{Bab \ref{ch:b3:rings} --- Ring dan Aritmetika}

\begin{solution}{exo:b3:rings:1}
(a) Pemasukan nilai $f \colon \Z[X] \to \Z[\iu]$, $P \mapsto P(\iu)$,
adalah morfisma ring surjektif ($a + bX \mapsto a + b\iu$). Kernelnya:
bagilah $P$ oleh polinomial \emph{monik} $X^2 + 1$ di $\Z[X]$: $P = (X^2 +
1)Q + (bX + a)$ dengan $a, b \in \Z$; lalu $P(\iu) = a + b\iu = 0$
jika dan hanya jika $a = b = 0$. Jadi $\ker f = (X^2+1)$ dan
\cref{thm:b3:rings:firstiso} menuntaskannya.

(b) Perhitungan yang sama dengan koefisien di $\R$: $\R[X]/(X^2+1)
\cong \C$ --- inilah \emph{konstruksi} $\C$ yang paling bersih.

(c) Polinomial $X^2 + X + 1$ tak berakar di $\mathbb F_2$ ($0, 1 \mapsto 1$),
sehingga, karena berderajat $2$, ia \omterm{def:b3:rings:divisibility}{tak tereduksi}: faktornya $\mathbb
F_4 = \mathbb F_2[X]/(X^2+X+1)$ adalah lapangan
(\cref{prop:b3:rings:primemaximal}; $(P)$ \omterm{def:b3:rings:primemaximal}{maksimal} di $K[X]$ bila
$P$ \omterm{def:b3:rings:divisibility}{tak tereduksi}, sebab $K[X]$ merupakan \omterm{def:b3:rings:pidufd}{DIU}: \omterm{def:b3:rings:ideal}{ideal} $(D) \supseteq
(P)$ berarti $D \mid P$). Keempat unsurnya adalah $0, 1, \omega,
\omega + 1$ dengan $\omega = \bar X$ dan $\omega^2 = \omega + 1$.
Tabel perkaliannya (untuk unsur tak nol):
\[
\omega \cdot \omega = \omega + 1, \qquad
\omega(\omega + 1) = \omega^2 + \omega = 1, \qquad
(\omega+1)^2 = \omega^2 + 1 = \omega .
\]
Unsur tak nolnya membentuk grup siklik berorde $3$ yang dibangun oleh
$\omega$.
\end{solution}

\begin{solution}{exo:b3:rings:2}
(a) Misalkan $p$ \omterm{def:b3:rings:divisibility}{tak tereduksi} di dalam sebuah \omterm{def:b3:rings:pidufd}{DFT} dan $p \mid ab$, katakanlah $ab =
pc$, dengan $a, b \neq 0$ (kalau tidak, trivial). Jika $a$ atau $b$ berupa unit,
maka $p$ membagi yang lain. Jika tidak, faktorkan $a$, $b$ dan $c$ atas
\omterm{def:b3:rings:divisibility}{unsur tak tereduksi}: kedua faktorisasi $ab$,
\[
(\text{faktor } a)(\text{faktor } b) = p \cdot
(\text{faktor } c),
\]
haruslah bersesuaian sampai pada urutan dan kesekawanan: jadi $p$ \omterm{def:b3:rings:divisibility}{sekawan} dengan salah satu
faktor \omterm{def:b3:rings:divisibility}{tak tereduksi} $a$ atau $b$, sehingga membaginya.

(b) Misalkan $A$ daerah integral berhingga dan $x \neq 0$. Pemetaan $y \mapsto
xy$ injektif ($xy = xy' \Rightarrow x(y - y') = 0 \Rightarrow
y = y'$), sehingga surjektif (karena $A$ berhingga): jadi $1 = xy$ untuk suatu $y$.

(c) Jika $\mathfrak p$ prima di dalam ring berhingga $A$, maka
$A/\mathfrak p$ daerah integral berhingga, sehingga lapangan menurut (b), jadi
$\mathfrak p$ \omterm{def:b3:rings:primemaximal}{maksimal} (\cref{prop:b3:rings:primemaximal}).
\end{solution}

\begin{solution}{exo:b3:rings:3}
Normanya: $N(3) = 9$, $N(1 \pm \iu\sqrt5) = 6$, $N(2 \pm \iu\sqrt5) =
9$. Persamaan $x^2 + 5y^2 = 2$ dan $x^2 + 5y^2 = 3$ tak punya
penyelesaian bulat, jadi tak ada unsur bernorma $2$ atau $3$. Faktorisasi sejati
dari $3$ akan memerlukan dua faktor bernorma $3$:
mustahil --- jadi $3$ \omterm{def:b3:rings:divisibility}{tak tereduksi}. Faktorisasi sejati dari $1
\pm \iu\sqrt5$ (bernorma $6$) akan memerlukan faktor bernorma $2, 3$:
mustahil juga. Demikian pula $2 \pm \iu\sqrt5$ (bernorma $9$: faktornya akan
bernorma $3$). Sekarang
\[
9 = 3 \cdot 3 = (2 + \iu\sqrt5)(2 - \iu\sqrt5),
\]
yaitu dua faktorisasi atas \omterm{def:b3:rings:divisibility}{unsur tak tereduksi}. Keduanya benar-benar
berbeda: unitnya hanya $\pm 1$ (bernorma $1$), dan $2 \pm \iu\sqrt5
\neq \pm 3$. Jadi ketunggalannya gagal --- sedangkan keberadaan
faktorisasi tetap berlaku di $\Z[\iu\sqrt5]$
(\cref{exo:b3:rings:10}(c)): ketakfaktorialan di sini murni
kegagalan ketunggalan. (Sejalan dengan itu, \cref{prop:b3:rings:primeirred}:
\omterm{def:b3:rings:divisibility}{unsur tak tereduksi} ini bukan \omterm{def:b3:rings:divisibility}{unsur prima}.)
\end{solution}

\begin{solution}{exo:b3:rings:4}
(a) Misalkan $a, b \in \Z[\iu]$, $b \neq 0$, dan $a/b = x + \iu y \in
\C$. Pilihlah bilangan bulat $m, n$ dengan $\abs{x - m} \leq \frac12$,
$\abs{y - n} \leq \frac12$, lalu tetapkan $q = m + \iu n$, $r = a - bq$.
Maka
\[
N(r) = N(b)\,\abs*{\tfrac ab - q}^2
\leq N(b)\Bigl(\tfrac14 + \tfrac14\Bigr) = \tfrac{N(b)}2 < N(b).
\]
Jadi $N$ merupakan fungsi \omterm{def:b3:rings:pidufd}{Euclid} ($N(r) < N(b)$ atau $r = 0$).

(b) Jika $uv = 1$ maka $N(u)N(v) = 1$ dengan $N(u) \in \N$: jadi $N(u) =
1$, yakni $x^2 + y^2 = 1$: $u \in \{\pm 1, \pm\iu\}$; sebaliknya
keempatnya memang unit.

(c) Untuk $\Z[\iu\sqrt2]$: pembulatan yang sama memberikan $\abs{a/b - q}^2
\leq \frac14 + \frac{2}4 = \frac34 < 1$: jadi \omterm{def:b3:rings:pidufd}{Euclid}; unitnya: $x^2 +
2y^2 = 1$ memberikan $\pm 1$. Untuk $\Z[\iu\sqrt5]$ batasnya menjadi
$\frac14 + \frac54 = \frac32 > 1$: argumen pembulatannya gagal ---
dan memang harus gagal, sebab $\Z[\iu\sqrt5]$ bahkan bukan \omterm{def:b3:rings:pidufd}{DFT}
(\cref{exo:b3:rings:3}), padahal \omterm{def:b3:rings:pidufd}{Euclid} akan mengakibatkan \omterm{def:b3:rings:pidufd}{DFT}
(\cref{thm:b3:rings:euclideanpid,thm:b3:rings:pidufd}).
\end{solution}

\begin{solution}{exo:b3:rings:5}
(a) Jika $x^n = 0$:
\[
(1 + x)\bigl(1 - x + x^2 - \dots + (-1)^{n-1}x^{n-1}\bigr)
= 1 + (-1)^{n-1}x^n = 1 .
\]
(b) Di daerah integral berlaku $\deg(PQ) = \deg P + \deg Q$; syarat $PQ = 1$ memaksa
$\deg P = \deg Q = 0$ dan $P, Q \in A^\times$: jadi $A[X]^\times =
A^\times$. Atas $\Z/4\Z$: $(1 + 2X)^2 = 1 + 4X + 4X^2 = 1$, sehingga $1
+ 2X$ merupakan unit berderajat $1$ (di sini $2$ nilpoten; bandingkan
dengan (a)).

(c) Dari $e^2 = e$ diperoleh $e(e - 1) = 0$, sehingga $e \in \{0, 1\}$ di
daerah integral. Jika $x^n = 0$ dengan $n \geq 1$ minimal dan $x \ne 0$, maka
$n \geq 2$ dan $x \cdot x^{n-1} = 0$ dengan kedua faktornya tak nol:
kontradiksi.
\end{solution}

\begin{solution}{exo:b3:rings:6}
(a) $K[X,Y]/(X,Y) \cong K$ (masukkan $(0,0)$): sebuah lapangan, jadi
$(X,Y)$ \omterm{def:b3:rings:primemaximal}{maksimal}. Jika $(X, Y) = (P)$: dari $P \mid X$ diperoleh (lihat derajatnya
dalam $Y$) $P \in K[X]$, lalu $P \mid Y$ memaksa $P \in K$; nilai $P =
0$ mustahil dan $P \in K^\times$ akan memberikan $(P) = K[X,Y]$,
bertentangan dengan kesejatiannya (sebab $K[X,Y]/(X,Y) \cong K \neq 0$). Jadi
$(X,Y)$ bukan \omterm{def:b3:rings:ideal}{ideal} utama.

(b) Pemasukan nilai $P(X, Y) \mapsto P(X, X^2)$ memetakan $K[X,Y]$ pada
$K[X]$; kernelnya adalah $(Y - X^2)$: dengan membagi oleh polinomial yang
monik dalam $Y$, yaitu $Y - X^2$, $P = (Y - X^2)Q + R(X)$, dan $P(X, X^2) =
R(X)$. Jadi $K[X,Y]/(Y - X^2) \cong K[X]$ --- ring koordinat sebuah
parabola, isomorfik dengan ring koordinat sebuah garis.

Pemasukan nilai $P(X, Y) \mapsto P(X, X^{-1})$ memetakan $K[X, Y]$ pada
ring $K[X, X^{-1}]$ yang beranggotakan polinomial Laurent. Kernelnya memuat
$(XY - 1)$; sebaliknya, modulo $XY - 1$ setiap kelas mempunyai
wakil $R = \sum_{n \geq 0} a_nX^n + \sum_{m \geq 1} b_m
Y^m$ (ganti setiap hasil kali $XY$ dengan $1$ berulang-ulang), dan $R(X,
X^{-1}) = \sum a_n X^n + \sum b_m X^{-m} = 0$ memaksa semua $a_n =
b_m = 0$. Karena itu $K[X, Y]/(XY - 1) \cong K[X, X^{-1}]$ --- yaitu
ring koordinat sebuah hiperbola: garis dengan satu titik dibuang.

(c) \omterm{def:b3:rings:ideal}{Ideal} $(Y - X^2)$ prima (sebab kuosiennya $K[X]$ daerah integral) tetapi tidak
\omterm{def:b3:rings:primemaximal}{maksimal} (sebab $K[X]$ bukan lapangan; secara konkret $(Y - X^2) \subsetneq
(Y - X^2,\, X) \subsetneq K[X,Y]$).
\end{solution}

\begin{solution}{exo:b3:rings:7}
\emph{$X^5 - 12X^3 + 36X - 12$}: \omterm{thm:b3:rings:criteria}{Eisenstein} pada $p = 3$ ($3 \mid
12, 36, 12$; $9 \nmid 12$; $3 \nmid 1$): jadi \omterm{def:b3:rings:divisibility}{tak tereduksi}. (Pada $p = 2$
\omterm{thm:b3:rings:criteria}{Eisenstein} gagal: $4 \mid 12$.)

\emph{$X^4 + X + 1$}: reduksikan modulo $2$. Tak berakar di $\mathbb F_2$;
satu-satunya polinomial kuadrat \omterm{def:b3:rings:divisibility}{tak tereduksi} atas $\mathbb F_2$ adalah $X^2 + X +
1$, dan $(X^2+X+1)^2 = X^4 + X^2 + 1 \neq X^4 + X + 1$. Jadi $X^4 +
X + 1$ \omterm{def:b3:rings:divisibility}{tak tereduksi} atas $\mathbb F_2$, sehingga juga atas $\Q$
(\cref{thm:b3:rings:criteria}(1); ia monik).

\emph{$X^4 + 4$}: terreduksi --- lewat identitas Sophie Germain,
$X^4 + 4 = (X^2 - 2X + 2)(X^2 + 2X + 2)$.

\emph{$X^4 + 1$}: geser, $(X+1)^4 + 1 = X^4 + 4X^3 + 6X^2 + 4X +
2$: \omterm{thm:b3:rings:criteria}{Eisenstein} pada $2$. Faktorisasi $X^4+1$ akan bergeser menjadi
faktorisasi $(X+1)^4 + 1$: jadi \omterm{def:b3:rings:divisibility}{tak tereduksi}.

\emph{$X^3 - X - 1$}: polinomial kubik terreduksi atas $\Q$ jika dan hanya jika ia berakar
rasional; akar rasional dari polinomial bulat yang monik pastilah
bilangan bulat pembagi suku konstannya (teorema akar rasional: jika
$(p/q)$ dalam bentuk paling sederhana merupakan akar, maka $q \mid 1$, $p \mid -1$), dan
$\pm 1$ bukan akar (nilainya $-1$ dan $-1$): jadi \omterm{def:b3:rings:divisibility}{tak tereduksi}.
\end{solution}

\begin{solution}{exo:b3:rings:8}
(a) Untuk $\gcd(m, n) = 1$, \cref{thm:b3:rings:crt} memberikan
isomorfisma ring $\Z/mn\Z \cong \Z/m\Z \times \Z/n\Z$. Unsur sebuah
ring hasil kali merupakan unit jika dan hanya jika kedua koordinatnya unit, sehingga
$(\Z/mn\Z)^\times \cong (\Z/m\Z)^\times \times (\Z/n\Z)^\times$
dan $\varphi(mn) = \varphi(m)\varphi(n)$. Untuk pangkat prima,
yang bukan unit di $\Z/p^k\Z$ adalah kelas kelipatan $p$:
jadi $\varphi(p^k) = p^k - p^{k-1}$. Karena itu
\[
\varphi(n) = \prod_i \bigl(p_i^{\alpha_i} -
p_i^{\alpha_i-1}\bigr) = n \prod_{p \mid n}\Bigl(1 -
\frac1p\Bigr).
\]

(b) $M = 7 \cdot 11 \cdot 13 = 1001$. Unsur idempotennya: $e_1 \equiv
(1, 0, 0)$: $143 = 11\cdot13 \equiv 3 \pmod 7$ dan $3 \cdot 5
\equiv 1$: jadi $e_1 = 143 \cdot 5 = 715$. Untuk $e_2$: $91 \equiv 3
\pmod{11}$, $3 \cdot 4 \equiv 1$: $e_2 = 91\cdot4 = 364$. Untuk $e_3$:
$77 \equiv -1 \pmod{13}$: $e_3 = 77 \cdot 12 = 924$. Maka
\[
x \equiv 2\,e_1 + 5\,e_2 + 1\,e_3 = 1430 + 1820 + 924 = 4174
\equiv 170 \pmod{1001},
\]
dan memang $170 = 24\cdot7 + 2 = 15\cdot11 + 5 = 13\cdot13 + 1$.
\end{solution}

\begin{solution}{exo:b3:rings:9}
(a) Jika $x^n = 0$ dan $y^m = 0$, penjabaran binomial
$(x+y)^{n+m}$ mempunyai setiap suku $x^iy^j$ dengan $i + j = n + m$, sehingga $i
\geq n$ atau $j \geq m$: setiap sukunya lenyap, dan $x + y$
nilpoten; lagi pula $(ax)^n = a^nx^n = 0$: jadi $\operatorname{Nil}(A)$ adalah
\omterm{def:b3:rings:ideal}{ideal}. Jika $\mathfrak p$ prima dan $x^n = 0 \in \mathfrak p$, maka
induksi pada $n$ memberikan $x \in \mathfrak p$ (sebab $x \cdot x^{n-1} \in
\mathfrak p$).

(b) Misalkan $a \notin \operatorname{Nil}(A)$ dan $S = \{a^n : n \geq
1\}$, sehingga $0 \notin S$. Himpunan $\mathcal E$ yang beranggotakan \omterm{def:b3:rings:ideal}{ideal} saling lepas
dengan $S$ memuat $(0)$ dan bersifat induktif (gabungan sebuah rantai
\omterm{def:b3:rings:ideal}{ideal} yang saling lepas dengan $S$ juga \omterm{def:b3:rings:ideal}{ideal} yang saling lepas dengan $S$): lema Zorn
menyediakan $\mathfrak p \in \mathcal E$ yang \omterm{def:b3:rings:primemaximal}{maksimal}. \omterm{def:b3:rings:ideal}{Ideal} $\mathfrak p$ ini
sejati (sebab $a \notin \mathfrak p$, karena $a \in S$). Keprimaannya: misalkan
$x, y \notin \mathfrak p$. Menurut kemaksimalan, $\mathfrak p + (x)$ dan
$\mathfrak p + (y)$ memotong $S$: $a^m \in \mathfrak p + (x)$, $a^n
\in \mathfrak p + (y)$. Dengan mengalikannya, $a^{m+n} \in \mathfrak p +
(xy)$. Jika $xy \in \mathfrak p$, maka $a^{m+n} \in \mathfrak p \cap
S$: mustahil. Jadi $xy \notin \mathfrak p$ --- inilah kontraposisi
keprimaan. Karena itu setiap unsur yang bukan nilpoten menghindari suatu \omterm{def:b3:rings:primemaximal}{ideal
prima}; bersama (a),
$\operatorname{Nil}(A) = \bigcap_{\mathfrak p} \mathfrak p$.
\end{solution}

\begin{solution}{exo:b3:rings:10}
(a) Rantai $\ker f \subseteq \ker f^2 \subseteq \cdots$
stabil (\cref{prop:b3:rings:noethacc}): $\ker f^n = \ker
f^{n+1}$ untuk suatu $n$. Misalkan $x \in \ker f$. Karena $f$, dan karenanya $f^n$,
surjektif, $x = f^n(y)$ untuk suatu $y$; lalu $f^{n+1}(y) = f(x) =
0$, sehingga $y \in \ker f^{n+1} = \ker f^n$, yakni $x = f^n(y) = 0$.

(b) Himpunan $I_n = \{f \in \mathcal C(\intcc01, \R) : f\restriction_{
\intcc0{1/n}} = 0\}$ adalah \omterm{def:b3:rings:ideal}{ideal}, dan $I_n \subseteq I_{n+1}$.
Inklusinya sejati: fungsi $x \mapsto \max\bigl(0, x -
\frac1{n+1}\bigr)$ lenyap pada $\intcc0{\frac1{n+1}}$ tetapi tidak pada
$\intcc0{\frac1n}$. Rantai naik sejati yang tak berujung
bertentangan dengan \cref{prop:b3:rings:noethacc}.

(c) Andaikan himpunan unsur tak nol bukan unit yang tak punya faktorisasi
atas \omterm{def:b3:rings:divisibility}{unsur tak tereduksi} bersifat tak kosong. Keluarga \omterm{def:b3:rings:ideal}{ideal}
$\{(a)\}$ yang bersesuaian mempunyai unsur \omterm{def:b3:rings:primemaximal}{maksimal} $(a)$
(\cref{prop:b3:rings:noethacc}). Unsur $a$ ini tidak
\omterm{def:b3:rings:divisibility}{tak tereduksi} (sebab \omterm{def:b3:rings:divisibility}{unsur tak tereduksi} adalah faktorisasi dirinya sendiri), jadi $a = bc$
dengan $b, c$ bukan unit; inklusi $(a) \subseteq (b)$ sejati (sebab $(a) =
(b)$ akan memberikan $b = ad$, $a = adc$, sehingga $dc = 1$: $c$ menjadi unit),
demikian pula $(a) \subsetneq (c)$. Menurut kemaksimalan, baik $b$ maupun $c$
terfaktorkan atas \omterm{def:b3:rings:divisibility}{unsur tak tereduksi}; menyambung faktornya memfaktorkan $a$:
kontradiksi. Diterapkan pada $\Z[\iu\sqrt5]$ --- yang \omterm{def:b3:rings:noetherian}{Noether} sebagai
kuosien dari $\Z[X]$ (\cref{thm:b3:rings:hilbert}, sebab $\Z[\iu\sqrt5]
\cong \Z[X]/(X^2+5)$) --- ini menunjukkan bahwa faktorisasi memang ada di sana;
\cref{exo:b3:rings:3} menunjukkan bahwa ketunggalanlah yang gagal.
\end{solution}

\begin{solution}{exo:b3:rings:11}
(a) \cref{exo:b3:rings:7}: geser lalu \omterm{thm:b3:rings:criteria}{Eisenstein} pada $2$.

(b) Modulo $2$: $X^4 + 1 = (X + 1)^4$. Sekarang misalkan $p$ ganjil. Semua
kuadrat membentuk subgrup berindeks $2$ di $(\Z/p\Z)^\times$:
morfisma $x \mapsto x^2$ berkernel $\{\pm 1\}$ (dua unsur, sebab $X^2
- 1$ berakar paling banyak $2$ di sebuah lapangan, dan $1 \neq -1$ untuk $p$
ganjil), sehingga petanya beranggota $\frac{p-1}2$ unsur. Akibatnya,
hasil kali dua bukan-kuadrat merupakan kuadrat (di dalam \omterm{thm:b3:groups:quotient}{grup kuosien}
berorde $2$ berlaku $\overline{xy} = \bar x\bar y$). Karena itu \emph{sekurang-kurangnya satu
dari $-1$, $2$, $-2$ merupakan kuadrat modulo $p$} (jika $-1$ dan $2$ bukan,
maka $-2 = (-1)\cdot 2$ merupakan kuadrat). Pada setiap kasus $X^4 + 1$ terfaktorkan modulo $p$:
\begin{itemize}
  \item $-1 = c^2$: \quad $X^4 + 1 = X^4 - c^2 = (X^2 - c)(X^2 +
        c)$;
  \item $2 = c^2$: \quad $(X^2 + cX + 1)(X^2 - cX + 1) = X^4 + (2
        - c^2)X^2 + 1 = X^4 + 1$;
  \item $-2 = c^2$: \quad $(X^2 + cX - 1)(X^2 - cX - 1) = X^4 -
        (c^2 + 2)X^2 + 1 = X^4 + 1$.
\end{itemize}
Jadi $X^4+1$ terreduksi modulo setiap bilangan prima, namun \omterm{def:b3:rings:divisibility}{tak tereduksi} atas
$\Q$: kriteria reduksi (\cref{thm:b3:rings:criteria}(1))
mendeteksi ketaktereduksian, tetapi kegagalannya tidak membuktikan apa pun.

(Alasan strukturalnya: $p^2 - 1 = (p-1)(p+1)$ merupakan hasil kali
dua bilangan genap berurutan, sehingga $8 \mid p^2 - 1$; grup siklik
$\mathbb F_{p^2}^\times$ (kesiklikannya dibuktikan pada
\cref{ch:b3:galois}) lalu memuat unsur $\zeta$ berorde
$8$, yaitu akar $X^4 + 1$; polinomial minimalnya atas $\mathbb
F_p$ membagi $X^4+1$ dan berderajat $\leq 2$ --- jadi $X^4+1$ tak pernah
dapat \omterm{def:b3:rings:divisibility}{tak tereduksi} modulo $p$.)
\end{solution}

\begin{solution}{exo:b3:rings:12}
(a) $(1-e)^2 = 1 - 2e + e^2 = 1 - e$. Pemetaan $\varphi(x) =
(ex, (1-e)x)$ bersifat aditif dan multiplikatif ke dalam hasil kali
kedua \omterm{def:b3:rings:ideal}{ideal} itu: $exey = e^2xy = e(xy)$, dan $Ae$ merupakan ring
komutatif bersatuan $e$ (sebab $e\cdot ex = ex$). Injektif:
dari $ex = 0$ dan $(1-e)x = 0$, penjumlahannya memberikan $x = 0$. Surjektif: $(ea,
(1-e)b)$ adalah peta dari $ea + (1-e)b$ (hitunglah kedua
komponennya dengan memakai $e(1-e) = 0$). Unitnya terpetakan secara benar pada
pasangan bertipe $(1, 0)$: $\varphi(1) = (e, 1-e)$, yaitu satuan
hasil kalinya.

(b) Di daerah integral, $e(e - 1) = 0$ memaksa $e \in \{0, 1\}$: jadi hanya
ada unsur idempoten trivial. Di $\Z/12\Z$, menyelesaikan $e^2 \equiv e$ memberikan $e
\in \{0, 1, 4, 9\}$. Pasangan tak trivialnya $\{4, 9\}$: $4 + 9 =
13 \equiv 1$, $4\cdot9 = 36 \equiv 0$, dan $\Z/12\Z\cdot4 =
\{0, 4, 8\} \cong \Z/3\Z$ (bersatuan $4$), $\Z/12\Z\cdot9 = \{0,
3, 6, 9\} \cong \Z/4\Z$ (bersatuan $9$): inilah pemecahan lewat teorema sisa Tiongkok,
$\Z/12\Z \cong \Z/4\Z\times\Z/3\Z$, dengan $9 \leftrightarrow
(1, 0)$ dan $4 \leftrightarrow (0, 1)$.

(c) Di bawah isomorfisma sisa Tiongkok $\Z/n\Z \cong \prod_i
\Z/p_i^{a_i}\Z$, unsur $e_i$ dengan kekongruenan yang
disebutkan bersesuaian dengan tupel yang bernilai $1$ pada slot $i$ dan
$0$ di tempat lain: itulah unsur idempoten elementer hasil kali tersebut.
Sebaliknya, keluarga lengkap unsur idempoten yang saling ortogonal
($e_ie_j = 0$ untuk $i \neq j$, $\sum e_i = 1$) merakit kembali
penguraian hasil kali itu lewat (a), secara induktif. Unsur idempoten bagi
ring adalah seperti proyeksi ortogonal bagi ruang Hilbert
(\cref{ch:b3:hilbert}): koordinat sebuah penguraian langsung
internal.
\end{solution}

\begin{solution}{pb:b3:rings:1}
\textbf{1.} $N(z) = z\bar z$, jadi $N(zw) = zw\overline{zw} = z\bar
z\, w \bar w = N(z)N(w)$. Jika $uv = 1$: $N(u)N(v) = 1$ di $\N$, sehingga
$N(u) = 1$, yakni $u \in \{\pm 1, \pm \iu\}$; keempatnya memang unit.
Brahmagupta: $(a^2+b^2)(c^2+d^2) = N\bigl((a + \iu b)(c + \iu
d)\bigr) = (ac - bd)^2 + (ad + bc)^2$.

\textbf{2.} Jika $z = ab$, maka $N(z) = N(a)N(b)$ berupa bilangan prima, sehingga
$N(a) = 1$ atau $N(b) = 1$: salah satu faktornya unit. Karena $N(z) > 1$,
$z$ bukan nol dan bukan unit: jadi \omterm{def:b3:rings:divisibility}{tak tereduksi} --- dan prima, sebab
$\Z[\iu]$ merupakan \omterm{def:b3:rings:pidufd}{DFT}
(\cref{thm:b3:rings:euclideanpid,thm:b3:rings:pidufd,%
lem:b3:rings:bezout}).

\textbf{3.} Unsur $\pi$ membagi $N(\pi) = \pi\bar\pi \geq 2$, sebuah
bilangan bulat; dengan memfaktorkan $N(\pi)$ atas bilangan prima dan memakai keprimaan
$\pi$, diperoleh $\pi \mid p$ untuk suatu bilangan prima $p$. Jika juga
$\pi \mid q \neq p$: Bézout di $\Z$ memberikan $1 = up + vq$, sehingga
$\pi \mid 1$ --- mustahil: jadi $p$ tunggal. Dari $p = \pi\gamma$:
$p^2 = N(p) = N(\pi)N(\gamma)$ dengan $N(\pi) \neq 1$, sehingga $N(\pi)
\in \{p, p^2\}$.

\textbf{4.} Misalkan $\pi$ prima Gauss yang membagi $p$, dengan $p =
\pi\gamma$. Jika $N(\pi) = p^2$: maka $N(\gamma) = 1$, sehingga $p$ \omterm{def:b3:rings:divisibility}{sekawan}
dengan $\pi$, yang juga prima Gauss; dan tak ada prima
Gauss bernorma $p$ (sebab jika $N(\rho) = p$ maka $\rho \mid
\rho\bar\rho = p$, dan keprimaan $p$ di $\Z[\iu]$ akan memaksa
$\rho$ \omterm{def:b3:rings:divisibility}{sekawan} dengan $p$, sehingga $N(\rho) = N(p) = p^2 \neq
p$).
Jika $N(\pi) = p$: dengan menulis $\pi = a + \iu b$, $p = \pi\bar\pi = a^2
+ b^2$.

\textbf{5.} Di dalam grup abelian $(\Z/p\Z)^\times$, pasangkan setiap
unsur dengan inversnya. Unsur yang menjadi inversnya sendiri adalah akar
$X^2 - 1$: tepat $\pm 1$ (sebuah lapangan berakar paling banyak dua).
Hasil kali semua unsurnya lalu $1 \cdot (-1) \cdot \prod
(\text{pasangan } k k^{-1}) = -1$: jadi $(p-1)! \equiv -1 \pmod p$. (Untuk
$p = 2$: $1! \equiv -1$.)

\textbf{6.} Tulis $(p-1)! = \prod_{k=1}^m k \cdot
\prod_{k=m+1}^{p-1}k$ dengan $m = \frac{p-1}2$. Pada hasil kali kedua
substitusikan $k = p - j$, $j = 1, \dots, m$: modulo $p$ berlaku
$\prod_{j=1}^m (p - j) \equiv (-1)^m m!$. Karena itu $-1 \equiv
(-1)^m (m!)^2$, yakni $(m!)^2 \equiv (-1)^{m+1} \pmod p$.

\textbf{7.} Jika $p \equiv 1 \pmod 4$, maka $m$ genap dan pertanyaan 6
memberikan $(m!)^2 \equiv -1$: yaitu akar kuadrat dari $-1$. Sebaliknya, jika
$x^2 \equiv -1 \pmod p$ (dengan $p$ ganjil), maka $x^4 = 1 \neq x^2$: jadi $x$
berorde $4$ di $(\Z/p\Z)^\times$, sehingga $4 \mid p - 1$ (Lagrange).
Dan untuk $p = 2$: $1^2 = 1 \equiv -1$. Kesimpulannya: $-1$ merupakan kuadrat
modulo $p$ jika dan hanya jika $p = 2$ atau $p \equiv 1 \pmod 4$.

\textbf{8.} Dengan $x^2 \equiv -1$: $p \mid x^2 + 1 = (x + \iu)(x -
\iu)$. Seandainya $p$ prima Gauss, ia akan membagi salah satu faktornya; padahal
$\frac xp \pm \frac \iu p \notin \Z[\iu]$. Jadi $p$ bukan
prima Gauss; menurut dikotomi (pertanyaan 4) --- $p$ yang tak prima
berarti cabang kedua --- $p = a^2 + b^2$.

\textbf{9.} Kuadrat bernilai $\equiv 0$ atau $1 \pmod 4$, sehingga $a^2 + b^2
\in \{0, 1, 2\} \pmod 4$: jadi bilangan prima $p \equiv 3 \pmod 4$ bukan
jumlah dua kuadrat. Menurut pertanyaan 4, cabang $N(\pi) = p$ (dengan $p =
a^2+b^2$) mustahil: jadi $p$ tetap menjadi prima Gauss.

\textbf{10.} $(1 + \iu)^2 = 2\iu$, sehingga $2 = -\iu(1 + \iu)^2$; dan
$N(1 + \iu) = 2$ berupa bilangan prima, jadi $1 + \iu$ merupakan prima Gauss
(pertanyaan 2).

\textbf{11.} Setiap prima Gauss membagi tepat satu bilangan
prima $p$ (pertanyaan 3); dengan mendaftarkannya menurut kasus: $p = 2$ memberikan
\omterm{def:b3:rings:divisibility}{sekawan} $1 + \iu$; $p \equiv 3 \pmod 4$ memberikan $p$ sendiri
(pertanyaan 9); $p \equiv 1 \pmod 4$ memberikan pasangan $\pi, \bar\pi$
bernorma $p$ (pertanyaan 4 dan 8). Pasangannya sungguhan: $\bar\pi \in
\{\pm\pi, \pm\iu\pi\}$ akan memaksa, dengan menulis $\pi = a + \iu b$,
entah $b = 0$, entah $a = 0$, entah $a = \pm b$, sehingga $p = a^2 + b^2 \in
\{a^2, 2a^2\}$ --- mustahil bagi bilangan prima ganjil. Periksa:
$5 = (2 + \iu)(2 - \iu)$, $N(2\pm\iu) = 5$; sedangkan $3$ prima bernorma
$9$.

\textbf{12.} Tulis $n = 2^{\alpha}\prod_i p_i^{\beta_i} \prod_j
q_j^{2\gamma_j}$ dengan $p_i \equiv 1$, $q_j \equiv 3 \pmod 4$.
Setiap faktornya merupakan jumlah dua kuadrat: $2 = 1^2 + 1^2$; $p_i = a^2
+ b^2$ (pertanyaan 8); $q_j^{2\gamma_j} = (q_j^{\gamma_j})^2 +
0^2$. Identitas Brahmagupta (pertanyaan 1) merambatkan
sifat itu ke hasil kalinya, yaitu $n$.

\textbf{13.} Misalkan $n = a^2 + b^2 = N(a + \iu b)$ dan $q \equiv 3
\pmod 4$ dengan $q \mid n$. Prima Gauss $q$ (pertanyaan 9) membagi
$(a + \iu b)(a - \iu b)$, sehingga membagi salah satu faktornya --- katakanlah $q
\mid a + \iu b$ (kasus lainnya serupa). Namun lalu $\frac{a +
\iu b}{q} = \frac aq + \iu \frac bq \in \Z[\iu]$ terbaca sebagai $q
\mid a$ \emph{dan} $q \mid b$ di $\Z$. Karena itu $q^2 \mid n$ dan
$\frac n{q^2} =
\bigl(\frac aq\bigr)^2 + \bigl(\frac bq\bigr)^2$. Dengan induksi kuat
pada $n$, eksponen $q$ di dalam $n/q^2$ genap; jadi eksponennya di dalam
$n$ pun genap.

\textbf{14.} \emph{Teorema (Fermat).} Bilangan bulat positif merupakan jumlah
dua kuadrat jika dan hanya jika setiap bilangan prima $\equiv 3 \pmod 4$
muncul padanya dengan eksponen genap. --- $2025 = 3^4 \cdot 5^2$:
eksponen $3$ genap, jadi ya ($2025 = 45^2 + 0^2 = 27^2 + 36^2$).
$2026 = 2 \cdot 1013$ dengan $1013 \equiv 1 \pmod 4$ prima: ya
($1013 = 22^2 + 23^2$, lalu Brahmagupta dengan $2 = 1^2+1^2$:
$2026 = (22 - 23)^2 + (22 + 23)^2 = 1^2 + 45^2$). Sedangkan $2027$ bilangan
prima $\equiv 3 \pmod 4$: tidak.

\textbf{15.} Misalkan $p = a^2 + b^2 = c^2 + d^2$ dengan bilangan bulat
positif, $p \equiv 1 \pmod 4$, dan $\pi$ prima Gauss dengan
$p = \pi\bar\pi$ (pertanyaan 4). Baik $a + \iu b$ maupun $c + \iu d$
bernorma $p$, sehingga keduanya prima Gauss (pertanyaan 2) yang membagi
$p = (a+\iu b)(a - \iu b)$; menurut ketunggalan faktorisasi, $c +
\iu d$ \omterm{def:b3:rings:divisibility}{sekawan} dengan $a + \iu b$ atau dengan $a - \iu b$:
\[
c + \iu d \in \{\pm(a \pm \iu b),\ \pm\iu(a \pm \iu b)\}
= \{\pm a \pm \iu b,\ \pm b \pm \iu a\}.
\]
Kepositifan $c, d$ menyisakan $c + \iu d \in \{a + \iu b, b + \iu
a\}$: jadi $\{c, d\} = \{a, b\}$.

\textbf{16.} Jika $n = a^2 - b^2 = (a-b)(a+b)$: kedua faktornya
berkeparitasan sama, sehingga $n$ ganjil (bila keduanya ganjil) atau habis dibagi
$4$ (bila keduanya genap) --- tak pernah $\equiv 2 \pmod 4$. Sebaliknya, untuk $n$
ganjil: $n = \bigl(\frac{n+1}2\bigr)^2 -
\bigl(\frac{n-1}2\bigr)^2$; dan untuk $n = 4m$: $n = (m+1)^2 - (m-1)^2$.
Jawabannya berupa syarat kekongruenan telanjang, dengan sebuah identitas
satu baris di baliknya: selisih kuadrat tak membawa kedalaman aritmetika,
dan kontrasnya dengan jumlah kuadrat itulah inti seluruh soal ini.

\textbf{17.} Pemetaan $(a, b) \mapsto z = a + \iu b$ adalah bijeksi
antara representasi dan $\{z : N(z) = n\}$. Faktorkan $z$ di dalam
\omterm{def:b3:rings:pidufd}{DFT} $\Z[\iu]$ dengan memakai klasifikasi (pertanyaan 11): sampai pada
sebuah unit, $z = (1+\iu)^{a'}\prod_j\pi_j^{s_j}\bar\pi_j^{t_j}
\prod_kq_k^{u_k}$, lalu ambil normanya ($N(1+\iu) = 2$,
$N(\pi_j) = N(\bar\pi_j) = p_j$, $N(q_k) = q_k^2$):
\[
n = 2^{a'}\prod_jp_j^{s_j + t_j}\prod_kq_k^{2u_k} .
\]
Dengan menyamakan eksponennya: $a' = a$, $s_j + t_j = b_j$, $2u_k = c_k$
--- \omterm{def:b3:groups:derived}{terselesaikan} jika dan hanya jika setiap $c_k$ genap, dan lalu $u_k = c_k/2$
terpaksa demikian sedangkan $s_j \in \intint0{b_j}$ bebas. Data
$(u, (s_j))$ yang berbeda memberikan $z$ yang tak \omterm{def:b3:rings:divisibility}{sekawan} dengan norma yang
\emph{sama}; unit $u \in \{\pm1, \pm\iu\}$ (4 pilihan) lalu
mencacah setiap kelas kesekawanan tanpa pengulangan (dua hasil kali yang sama
akan melanggar ketunggalan faktorisasi ---
$\pi_j$ dan $\bar\pi_j$ tidak \omterm{def:b3:rings:divisibility}{sekawan} sebab $p_j =
\pi_j\bar\pi_j$ tidak bercabang). Totalnya: $r_2(n) =
4\prod_j(b_j + 1)$, dan $0$ jika ada $c_k$ yang ganjil.

\textbf{18.} Kesamaan $\chi(dd') = \chi(d)\chi(d')$ diperiksa modulo $4$
(ganjil $\times$ ganjil mencakup keempat kasus tandanya; sedangkan bilangan genap
memberikan $0 = 0$). Untuk $m, n$ yang relatif prima, pembagi $mn$
secara tunggal berbentuk $d = d_1d_2$ dengan $d_1 \mid m$, $d_2 \mid n$:
$\sum_{d \mid mn}\chi(d) = \bigl(\sum_{d_1\mid
m}\chi(d_1)\bigr)\bigl(\sum_{d_2\mid n}\chi(d_2)\bigr)$: jadi
multiplikatif. Pada pangkat prima: di $2^a$ hanya $d = 1$ yang ganjil,
sehingga jumlahnya $= 1$. Di $p^b$ dengan $p \equiv 1$: semua $\chi(p^i) = 1$,
jumlahnya $= b + 1$. Di $q^c$ dengan $q \equiv 3$: $\chi(q^i) =
(-1)^i$, sehingga jumlah bergantinya $= 1$ (bila $c$ genap) atau $0$ (bila $c$ ganjil).

\textbf{19.} Kedua fungsi multiplikatif
$\frac14r_2$ (pertanyaan 17) dan $\sum_{d\mid n}\chi(d)$
(pertanyaan 18) bersesuaian pada semua pangkat prima --- $1$ di $2^a$;
$b + 1$ di $p^b$; $\mathbf 1_{c\ \mathrm{even}}$ di $q^c$ ---
sehingga bersesuaian di mana-mana: itulah rumus Jacobi, dengan $\sum_{d\mid
n}\chi(d) = d_1(n) - d_3(n)$ setelah pembaginya dipilah. Pemeriksaan:
$r_2(3) = 0 = 4(1 - 1)$; $r_2(5) = 8 = 4(2 - 0)$
(yaitu $(\pm1,\pm2), (\pm2,\pm1)$); $r_2(9) = 4 = 4(2 - 1)$
(pembaginya $1, 9 \equiv 1$; $3 \equiv 3$; representasinya
$(\pm3, 0), (0, \pm3)$); $r_2(25) = 12 = 4(3 - 0)$. Untuk $65 =
5\cdot13$: $r_2 = 4\cdot2\cdot2 = 16$, dari $65 = 1 + 64 = 16
+ 49$: yaitu keenam belas pasangan $(\pm1, \pm8), (\pm8, \pm1), (\pm4,
\pm7), (\pm7, \pm4)$.

\textbf{20.} Jumlah $\sum_{n \leq x}r_2(n)$ mencacah pasangan $(a, b)$
dengan $0 < a^2 + b^2 \leq x$, yakni titik kisi pada
cakram tertutup $D_{\sqrt x}$ dikurangi titik asalnya. Kaitkan pada setiap
titik kisi $P$ sebuah persegi satuan $P + \intco01^2$: persegi-persegi itu
memetak-metak bidang. Setiap persegi yang terkait pada titik di
$D_{\sqrt x}$ terletak di $D_{\sqrt x + \sqrt2}$, dan setiap persegi yang
memotong $D_{\sqrt x - \sqrt 2}$ terkait pada titik di
$D_{\sqrt x}$ (sebab perseginya berdiameter $\sqrt 2$): dengan membandingkan
luasnya,
areas,
\[
\pi(\sqrt x - \sqrt2)^2 \leq \#\{\text{titik kisi di }
D_{\sqrt x}\} \leq \pi(\sqrt x + \sqrt 2)^2,
\]
dan kedua batasnya sama dengan $\pi x + O(\sqrt x)$. Membuang titik asalnya
tak mengubah apa pun pada ketelitian ini.

\textbf{21.} Menurut Jacobi (pertanyaan 19) dan setelah urutan
penjumlahannya ditukar ($n = dm$):
\[
\frac14\sum_{n\leq x}r_2(n) = \sum_{n \leq x}\sum_{d \mid
n}\chi(d) = \sum_{d \leq x}\chi(d)\,\#\{m : dm \leq x\}
= \sum_{d\leq x}\chi(d)\Bigl\lfloor\frac xd\Bigr\rfloor,
\]
yang sama dengan $\frac{\pi x}4 + O(\sqrt x)$ menurut pertanyaan 20. Lepaskan
fungsi lantainya: $\lfloor x/d\rfloor = x/d + O(1)$, tetapi menjumlahkan
$O(1)$ atas $d \leq x$ terlalu kasar; sebagai gantinya pakailah bahwa
jumlah parsial $\chi$ terbatas ($0, 1, 1, 0$ secara siklik),
sehingga lewat penjumlahan Abel $\sum_{d\leq x}\chi(d)\{x/d\}$, yang
sukunya kita kelompokkan berpasangan menurut $d \equiv 1, 3$, bernilai $O(\sqrt x)$ ---
atau, lebih sederhana lagi: pecahlah pada $\sqrt x$. Untuk $d \leq
\sqrt x$, ganti $\lfloor x/d\rfloor$ dengan $x/d + O(1)$: galatnya
$O(\sqrt x)$. Untuk $d > \sqrt x$, nilai $\lfloor x/d\rfloor$ mengambil
setiap nilai $v < \sqrt x$ pada satu selang $d$ yang berurutan,
dan di situ jumlah $\chi$ bernilai $O(1)$: jadi galat totalnya $O(\sqrt x)$ setelah
dijumlahkan atas $\leq \sqrt x$ nilai $v$, sedangkan
$\sum_{d > \sqrt x}\chi(d)\frac xd = O(\sqrt x)$ menurut
ekor deret berganti tanda ($x\sum_{d>\sqrt x}\chi(d)/d =
x\,O(1/\sqrt x)$). Karena itu
\[
x\sum_{d \leq x}\frac{\chi(d)}d = \frac{\pi x}4 + O(\sqrt x),
\qquad\text{yakni}\qquad
\sum_{d\leq x}\frac{\chi(d)}d = \frac\pi4 +
O\Bigl(\frac1{\sqrt x}\Bigr),
\]
dan dengan membiarkan $x \to \infty$: $1 - \frac13 + \frac15 - \dots =
\frac\pi4$.

\textbf{22.} Seandainya $n \equiv 3 \pmod4$ sama dengan $a^2 + b^2$: kuadrat
bernilai $\equiv 0, 1 \pmod 4$, dan $a^2 + b^2 \in \{0, 1, 2\}$
modulo $4$ --- mustahil. (Kriteria pertanyaan 17 mengatakan hal yang
sama: $n \equiv 3 \pmod 4$ memaksa suatu bilangan prima $\equiv 3$ berpangkat
ganjil.) Jadi bilangan bulat yang terwakilkan menghindari satu kelas sisa
penuh: kepadatannya $\leq \frac34$. Rata-rata $\pi$ dari $r_2$
memusat pada segelintir bilangan bulat: $n = \prod_{j\leq k}p_j$
(dengan bilangan prima berbeda $\equiv 1 \bmod 4$) mempunyai $r_2(n) = 4\cdot2^k$
representasi --- tak terbatas banyaknya --- sehingga himpunan $n$ yang
jarang dapat membawa seluruh rata-ratanya, persis seperti nilai harapan
sebuah lotre berdampingan dengan kekalahan yang hampir pasti. Hasil Landau,
$\#\{n \leq x \text{ terwakilkan}\} \sim
Cx/\sqrt{\log x}$, membenarkannya: kepadatan $0$, rata-rata $\pi$.

\textbf{23.} \emph{Keperluan.} Misalkan $n = a^2 + b^2$ dengan
$\gcd(a, b) = 1$. Seandainya bilangan prima $q \equiv 3 \pmod 4$ membagi $n$,
pertanyaan 13 menunjukkan $q \mid a$ dan $q \mid b$: kontradiksi. Seandainya
$4 \mid n$: kuadrat bernilai $\equiv 0, 1 \pmod 4$, jadi $a^2 + b^2
\equiv 0 \pmod 4$ memaksa $a^2 \equiv b^2 \equiv 0$, yakni $a, b$
sama-sama genap: kontradiksi. \emph{Kecukupan.} Tulis $n =
2^{\alpha}\prod_jp_j^{b_j}$ dengan $\alpha \leq 1$ dan $p_j
\equiv 1 \pmod 4$, lalu tetapkan $z = (1+\iu)^{\alpha}\prod_j
\pi_j^{b_j} = a + \iu b$, yang bernorma $n$. Andaikan bilangan prima $t$
membagi $\gcd(a, b)$; maka $t \mid z$ di $\Z[\iu]$. Jika $t
\equiv 3 \pmod 4$: $t \mid N(z) = n$, dan itu sudah dikecualikan. Jika $t \equiv 1
\pmod 4$: $t = \pi_t\bar\pi_t$, sehingga $\bar\pi_t \mid z$; padahal
faktorisasi $z$ tak memuat prima \omterm{def:b3:rings:divisibility}{sekawan} mana pun (sebab $\pi_j$ dan
$\bar\pi_j$ tidak \omterm{def:b3:rings:divisibility}{sekawan}, pertanyaan 17), dan itu bertentangan dengan
ketunggalan faktorisasi. Jika $t = 2 = -\iu(1+\iu)^2$: maka
$(1+\iu)^2 \mid z$, yang memaksa $\alpha \geq 2$, juga dikecualikan. Karena itu
$\gcd(a, b) = 1$: representasinya primitif.

\textbf{24.} Bilangan $a$ ganjil (sebab $\gcd(a, b) = 1$ dan $b$ genap), sehingga $c^2
= a^2 + b^2$ ganjil dan $c$ pun ganjil. Misalkan $\delta$ sebuah pembagi
prima Gauss persekutuan dari $a + \iu b$ dan $a - \iu b$: ia
membagi jumlahnya $2a$ dan selisihnya $2\iu b$, sehingga membagi
$2a$ dan $2b$; hubungan Bézout $ua + vb = 1$ lalu memberikan
$\delta \mid 2$, jadi $\delta$ \omterm{def:b3:rings:divisibility}{sekawan} dengan $1 + \iu$ dan
$N(\delta) = 2$ membagi $N(a + \iu b) = c^2$, padahal itu ganjil:
kontradiksi. Jadi $a + \iu b$ dan $a - \iu b$ relatif prima dengan
hasil kali $c^2$; di dalam \omterm{def:b3:rings:pidufd}{DFT} $\Z[\iu]$, setiap prima Gauss dari
$c^2$ muncul dengan eksponen genap dan seluruhnya masuk ke salah satu
dari kedua faktor yang relatif prima itu, sehingga $a + \iu b = u(m + \iu
n)^2 = u\bigl(m^2 - n^2 + 2\iu mn\bigr)$ dengan $u$ sebuah unit. Pilihan
$u = \pm\iu$ membuat bagian realnya $\mp 2mn$ menjadi genap ---
mustahil, sebab $a$ ganjil. Pilihan $u = \pm1$ memberikan, setelah
tanda $m, n$ disesuaikan dan namanya ditukar agar
semuanya positif, $a = m^2 - n^2$, $b = 2mn$ dengan $m > n
\geq 1$; dan $c^2 = N(m + \iu n)^2$ memberikan $c = m^2 + n^2$. Pembagi
persekutuan $m$ dan $n$ akan membagi $a$ dan $b$: jadi $\gcd
(m, n) = 1$; sedangkan $m \equiv n \pmod 2$ akan membuat $a$ genap: jadi
keparitasannya berlawanan. Pemeriksaan: $(m, n) = (2, 1)$ memberikan $(3, 4,
5)$; $(m, n) = (5, 2)$ memberikan $(25 - 4, 20, 25 + 4) = (21, 20,
29)$, dan $441 + 400 = 841 = 29^2$.

\textbf{25.} Rumus Jacobi $r_2(n) = 4(d_1(n) - d_3(n))$
memberikan, untuk $n = 1, \dots, 25$:
\[
4,\ 4,\ 0,\ 4,\ 8,\ 0,\ 0,\ 4,\ 4,\ 8,\ 0,\ 0,\ 8,\ 0,\ 0,\
4,\ 8,\ 4,\ 0,\ 8,\ 0,\ 0,\ 0,\ 0,\ 12,
\]
yang tak nol tepat pada $n = 1, 2, 4, 5, 8, 9, 10, 13, 16, 17, 18,
20, 25$ (misalnya $r_2(15) = 0$: pembaginya $1, 5 \equiv 1$
dan $3, 15 \equiv 3$ saling mengimbangi; sedangkan $r_2(20) = 8$: pembaginya $1, 5
\equiv 1$, tak ada yang $\equiv 3$). Totalnya $4 + 4 + 4 + 8 + 4 +
4 + 8 + 8 + 4 + 8 + 4 + 8 + 12 = 80$. Dari sisi pembaginya:
$d \leq 25$ yang ganjil menyumbang
\[
25 - 8 + 5 - 3 + 2 - 2 + 1 - 1 + 1 - 1 + 1 - 1 + 1 = 20,
\]
yang terbaca sebagai $\chi(d)\lfloor 25/d\rfloor$ untuk $d = 1, 3, 5, \dots,
25$; dan $4 \cdot 20 = 80$, seperti diramalkan identitas pertanyaan 21.
Titik kisi pada cakram tertutup berjari-jari $5$: ada
$80$ titik dengan $1 \leq a^2 + b^2 \leq 25$ ditambah titik
asalnya, yakni $81$; sedangkan $\pi x = 25\pi \approx 78.54$, sehingga
galatnya sekitar $2.46$, masih nyaman di dalam pita $O(\sqrt x)$
pada pertanyaan 20 ($\sqrt x = 5$).

\end{solution}
